% Full-page figure: schematic decay of singular values for several
% classes of engineering matrix. Log-y axis. The figure carries a
% pedagogical message: the engineer reads "effective rank" off the
% curve.
\begin{figure}[p]
\centering
\begin{tikzpicture}
\begin{axis}[
  width=14cm,
  height=18cm,
  xlabel={index $k$},
  ylabel={singular value $\sigma_{k}/\sigma_{1}$},
  xmin=1, xmax=200,
  ymin=1e-6, ymax=2,
  ymode=log,
  grid=major,
  grid style={draw=enggray, very thin, opacity=0.4},
  legend pos=south west,
  legend cell align=left,
  legend style={font=\small, fill=white, draw=engblack, thick},
  tick label style={font=\footnotesize},
  label style={font=\small},
  title={Singular-value decay across matrix classes},
  title style={font=\bfseries},
]

% Class 1: low-rank-plus-noise: first 10 are O(1), then plateau at 1e-4
\addplot[engred, very thick, mark=none, samples=200, domain=1:200]
  {(x <= 10) * (1 - (x - 1) * 0.08) + (x > 10) * 0.0002};
\addlegendentry{low-rank + noise (effective rank $\approx 10$)}

% Class 2: natural image, ~1/k decay
\addplot[engblue, very thick, dashed, mark=none, samples=200, domain=1:200]
  {1/x};
\addlegendentry{natural image ($\sigma_{k} \propto 1/k$)}

% Class 3: random Gaussian matrix, slow decay
\addplot[engamber, very thick, mark=none, samples=200, domain=1:200]
  {1 - 0.6 * (x - 1) / 200};
\addlegendentry{random Gaussian (no spectral gap)}

% Class 4: covariance from spiked model, sharp drop after k=5
\addplot[engblack, very thick, dotted, mark=none, samples=200, domain=1:200]
  {(x <= 5) * (1 - (x - 1) * 0.15) + (x > 5) * (0.5 / (x - 4))};
\addlegendentry{spiked covariance (5 strong factors)}

% Reference horizontal lines: noise floor at 1e-3, threshold at 1e-2
\addplot[gray, thin, dashed, domain=1:200] {1e-2};
\node[font=\footnotesize, anchor=west, gray] at (axis cs:160, 1.4e-2)
  {1\% threshold};

\addplot[gray, thin, dashed, domain=1:200] {1e-4};
\node[font=\footnotesize, anchor=west, gray] at (axis cs:160, 1.4e-4)
  {numerical noise floor (single precision)};

\end{axis}
\end{tikzpicture}
\caption[Singular-value decay schematics]{Singular value decay
governs the cost-quality trade in low-rank approximation. The
horizontal axis is the singular-value index $k$; the vertical axis is
$\sigma_{k}/\sigma_{1}$ on a logarithmic scale. Four schematic
classes: a low-rank-plus-noise matrix has a sharp drop at the
effective rank and a plateau at the noise level (red); a natural
image's singular values decay roughly as $1/k$ over the first few
hundred indices (blue dashed); a random Gaussian matrix has no
spectral gap and singular values nearly uniform until the matrix
boundary (amber); a spiked-covariance model has a sharp leading
spectrum followed by a slow tail (black dotted). The horizontal lines
mark the $1\%$ relative-singular-value threshold (the engineer's
rule-of-thumb truncation level for many applications) and the
numerical noise floor in single-precision arithmetic (below which
singular values cannot be trusted as signal). The engineer reading
this plot can pick a truncation rank $k$ where $\sigma_{k}/\sigma_{1}$
crosses the application's acceptable error and verify that the data
class is one of the favourable ones. A flat curve like the amber
random-Gaussian schematic says that truncation will cost accuracy and
the matrix is genuinely full rank; a curve like the red low-rank
schematic says that nearly all the action is in the leading
components and aggressive truncation is safe.}
\label{fig:vol02:ch10:singular-value-decay}
\end{figure}
